
\chapter{Подписки и сохранённые реквизиты}\label{ch:subscriptions}
Подписка начинается с~первого согласия на хранение реквизитов. Для
продукта январский платёж и февральское автосписание продолжают
одну услугу, но для эмитента это две разные авторизационные
ситуации. Если связь между ними обрывается, эмитент теряет контекст,
схема считает операцию подозрительной, а~клиент получает отказ там,
где платёж должен был пройти.

\section{Сохранённые реквизиты и SCF}\label{sec:subscriptions-scf}

До появления единых правил продавцы решали повторные списания каждый
по-своему: кто-то хранил PAN, кто-то полагался на токенизацию, кто-то
каждый месяц отправлял обычный запрос и~надеялся, что эмитент угадает
контекст. Главное эмитенту не сообщали: продолжение ли это
согласованной цепочки или новая операция без явного участия клиента.

Отсутствие контекста делало \enquote{тихую} MIT неотличимой от обычного
CNP-запроса. Продавец мог начать регулярные списания без явного
согласия клиента или продолжить их после его отмены, и~эмитент
не имел технической возможности это различить. SCF закрыл этот
разрыв: контекст первого согласия стал обязательным полем
авторизационного сообщения, а~не опциональным пожеланием.

Visa и Mastercard ввели Stored Credential Framework (SCF) -- набор
правил, по которым схема различает первое согласие клиента
и последующие списания, делая контекст первого согласия видимым
в~авторизационном сообщении~\cite{visa-scf,mc-tpr}. Карта остаётся
у~продавца как сохранённый реквизит, но каждая операция должна
ясно показывать связь с~исходным согласием клиента.

В~авторизационном сообщении эта связь выражается несколькими
признаками: первичным или последующим использованием сохранённых
реквизитов, типом сценария и~ссылкой на исходную авторизацию, связывающей
всю цепочку~\cite{mc-tpr,cybs-initial-transactions}.
На уровне авторизационного сообщения SCF сводится к~нескольким
конкретным индикаторам. Visa использует Stored Credential
Transaction Indicator со значением \texttt{initial} при первой
CIT и \texttt{subsequent} при последующей MIT. Рядом передаётся
тип инициатора (cardholder-initiated или merchant-initiated)
и~причина MIT -- recurring, installment или
unscheduled~\cite{visa-scf,cybs-initial-transactions}. Mastercard
кодирует ту же информацию через subelements DE~48 (POS Entry
Mode и MIT Framework Indicator; точные subfields сверяются
по действующей редакции TPR) и~требует передавать Trace~ID
исходной CIT в~каждой последующей MIT~\cite{mc-tpr}.

Пример: клиент оформляет месячную подписку. Первый платёж идёт
как CIT с indicator = \texttt{initial} и
получает от схемы идентификатор транзакции (NTID у Visa, Trace~ID
у Mastercard). Через 30 дней продавец отправляет MIT с
indicator = \texttt{subsequent}, типом \texttt{recurring} и
ссылкой на NTID первой CIT. Если продавец не передаст ссылку
на исходную авторизацию или укажет неверный тип, эмитент увидит
запрос без контекста и вправе применить стандартную CNP-политику,
вплоть до отказа. Схема при аудите может квалифицировать такую
операцию как нарушение правил SCF~\cite{mc-tpr,visa-core-rules}.

\begin{figure}[tbp]
  \centering
  \includefiguresvg[width=.8\linewidth]{assets/figures/ch22-subscriptions/scf-lifecycle}
  \caption{Жизненный цикл SCF: CIT создаёт согласие, MIT продолжает цепочку.}\label{fig:scf-lifecycle}
\end{figure}

\section{CIT и MIT}\label{sec:subscriptions-cit-mit}

Подписочный процесс ломается между двумя шагами. Клиент
оформил услугу и согласился на сохранение реквизитов, а~через
неделю или месяц продавец пытается списать очередной платёж уже
без его участия. На этом стыке и обрывается контекст.

Транзакция, инициированная клиентом (customer-initiated transaction,
CIT) -- первый шаг. Клиент сам нажимает кнопку оплаты,
вводит данные карты, проходит аутентификацию, если она требуется
правилами рынка. Так выглядит первый платёж за подписку, оформление
Card-on-File или покупка, после которой продавец получает право
сохранить реквизиты. В~этот момент клиент платит, и~система фиксирует
исходную точку цепочки, получая связующий идентификатор исходной
авторизации. В~документации сети или шлюза он называется
transaction identifier, network transaction ID,
\texttt{previousTransactionID} или Trace ID~\cite{mc-tpr,cybs-initial-transactions,cybs-mc-subscription-mit}.

Согласие клиента возникает в CIT, но сохранение реквизитов
не означает право на любые будущие списания. Идентификатор нужно
хранить и передавать в~каждой последующей операции, иначе связь
между первым согласием и следующими списаниями исчезнет из
авторизационного контекста~\cite{mc-tpr,cybs-initial-transactions}.

Транзакция, инициированная продавцом (merchant-initiated transaction,
MIT) -- второй шаг. Это списание без активного участия клиента:
ежемесячная подписка, доначисление по уже оказанной услуге или
платёж по ранее согласованной схеме. Запрос инициирует продавец,
но обязан показать эмитенту, что операция продолжает существующую
цепочку. Поэтому MIT всегда живёт рядом со SCF, схемными признаками
типа операции и передачей идентификатора исходной
авторизации~\cite{mc-tpr,cybs-mc-subscription-mit}.

Ограничение жёсткое, и~MIT без правильного контекста расценивается
как новая и подозрительная транзакция. В~европейском трафике это
особенно заметно, потому что при удалённом установлении MIT-мандата
сильная аутентификация клиента (strong customer authentication, SCA)
применяется на этапе создания мандата, а последующие MIT должны быть
корректно помечены и~связаны с~исходной
авторизацией~\cite{eba-mit-qa,mc-tpr}. Если инициатор и контекст
описаны неверно, эмитент отклоняет платёж даже тогда, когда
бизнес-логика продавца считает его штатным.

\section{Идентификатор исходной авторизации и цепочка согласия}\label{sec:subscriptions-ntid}

Токен отвечает на вопрос, каким реквизитом платить, а идентификатор
исходной авторизации отвечает на вопрос, с~каким ранее данным
согласием связано текущее автосписание. Это разные
сущности, и смешивать их опасно. Токен можно обновить при перевыпуске
карты, а идентификатор цепочки должен храниться отдельно от самих
реквизитов. Mastercard в~ЕЭЗ, Великобритании и Гибралтаре прямо требует
в~последующих запросах авторизации по регулярным платежам передавать
уникальный Trace ID из ответа на исходную авторизацию, а~эмитент должен
уметь по нему восстановить исходную транзакцию~\cite{mc-tpr}. В
практической документации Visa/Cybersource та же идея сформулирована
ещё проще. После начальной операции нужно сохранить идентификатор
транзакции, transaction identifier, потому что он понадобится для
последующих операций~\cite{cybs-initial-transactions}.

Терминология различается между схемами. Visa называет
этот идентификатор Network Transaction ID (NTID) и
возвращает его в~ответе на авторизацию; Mastercard использует
термин Trace~ID и~требует его передачу в DE~48 последующих
авторизаций~\cite{mc-tpr,cybs-initial-transactions}. У
Cybersource и других шлюзов поле может называться
\texttt{previousTransactionID} или
\texttt{networkTransactionID} -- разные имена для одной
сущности.

Хранить идентификатор нужно на уровне подписки, а~не на уровне
транзакции. Каждый последующий MIT ссылается на одну и~ту же
исходную CIT, поэтому идентификатор живёт столько же, сколько
сама подписка. Если система хранит его только в~записи
последней транзакции, при любом сбое записи цепочка обрывается.
На практике это означает отдельное поле в~модели подписки,
которое заполняется при первой CIT и~не перезаписывается
последующими ответами~\cite{cybs-mc-subscription-mit}.

Типичный сбой возникает при миграции между PSP. Продавец переносит
токены и~историю подписок, но забывает перенести идентификатор
исходной авторизации первой CIT. На следующем цикле MIT приходит уже
без этого якоря, и~продавец теряет связующий идентификатор
между первым согласием и последующим списанием~\cite{mc-tpr,cybs-mc-subscription-mit}.

\begin{figure}[tbp]
  \centering
  \includefiguresvg[width=.8\linewidth]{assets/figures/ch22-subscriptions/ntid-chain}
  \caption{Цепочка NTID: от первой CIT через MIT к~точке обрыва при миграции PSP.}\label{fig:ntid-chain}
\end{figure}

Идентификатор хранится как обязательный атрибут подписки,
передаётся во всех последующих MIT и~не разрывается при смене
процессора. Это требует той же аккуратности, что и~работа с
токенами. Если исходная CIT оформлена некорректно, последующий
косметический ремонт цепочку согласия восстановить не сможет.

\section{Мягкий отказ и повторная попытка}\label{sec:subscriptions-soft-decline}

Когда MIT получает отказ, который не закрывает сценарий окончательно,
допустимых действий несколько: новый шаг аутентификации, повтор
в~разрешённом окне, обработка кода Merchant Advice Code
или уведомление клиента~\cite{visa-core-rules,mc-tpr}.

Если отказ требует аутентификации, клиента возвращают в~новый
подтверждённый шаг, а~не пытаются задним числом \enquote{додавить} ту
же MIT. Европейское регулирование и схемные правила здесь сходятся:
при удалённом установлении MIT-мандата SCA применяется на этапе
создания мандата, а~не после произвольной серии
повторов~\cite{eba-mit-qa,mc-tpr}.

Если же речь не об аутентификации, а~об отказе по регулярному списанию,
то действуют уже конкретные схемные процедуры. Visa для подписочных
сценариев устанавливает отдельные требования к~уведомлению держателя
карты -- например, при окончании пробного или промо-периода продавец
обязан предупредить клиента не менее чем за 7 календарных дней
до первого полного списания; для повторной отправки отклонённой
MIT-операции по разрешающим кодам отказа (\textit{decline}) Visa
ограничивает количество попыток (до 15 в~30-дневном
окне)~\cite{visa-core-rules}. Mastercard для отказов по
регулярным платежам ограничивает повторную отправку (до 10~попыток
за 24~часа, до 35~попыток за 30~дней) и~требует учитывать Merchant
Advice Code (MAC), который определяет допустимый интервал
и~саму возможность повтора~\cite{mc-tpr}. Разница между лимитами
MC (35/30 дней) и Visa (15/30 дней) отражает разные методы подсчёта.
Оба лимита выставлены в~зоне, где совокупный уровень успеха повторов
перестаёт расти, а~нагрузка на авторизационную инфраструктуру
продолжает расти.

MAC определяет и разрешение на повтор, и конкретное действие после
отказа. Коды MAC~24--30 задают минимальное время ожидания перед
повтором: MAC~24 -- 1~час, MAC~25 -- 24~часа, MAC~26 -- 2~дня,
MAC~27 -- 4~дня, MAC~28 -- 6~дней, MAC~29 -- 8~дней, MAC~30 -- 10~дней.
Эскалирующее окно устроено не произвольно. Код~24 рассчитан на
временное ограничение баланса, которое снимется к~следующему
расчётному дню, а~старшие коды -- на ситуацию, когда держатель ждёт
пополнения в~следующем расчётном цикле. Схема \enquote{обучает} ТСП
откладывать повтор до момента, когда вероятность одобрения
повышается. MAC~03
(\textit{Do not try again}) полностью запрещает повторную
отправку по этим реквизитам, а MAC~01
(\textit{New account information available}) сигнализирует, что
реквизиты устарели и~продавец должен получить обновлённые данные
через ABU или от клиента~\cite{mc-tpr}. Игнорирование MAC
приводит к~штрафному трафику, и~Mastercard отслеживает повторные
попытки после запрета и~включает ТСП в~мониторинговую
программу.

Visa использует схожую логику через Response Code и Transaction
Advisor Code. Отказ с~кодом, требующим аутентификации, возвращает
клиента в 3DS-цикл, а~не в~повторную MIT. Отказ с~кодом,
допускающим повтор, ограничен 15~попытками за 30~дней для
одного ТСП и~набора реквизитов~\cite{visa-core-rules}.
Дальнейшие попытки Visa расценивает как нарушение правил.

Биллинговая система прогоняет каждый отказ через классификатор:
допустим ли повтор, через какое время, нужно ли вернуть клиента
в~интерактивный сценарий. Без такого классификатора биллинг
либо повторяет слишком часто и~попадает под мониторинг схемы,
либо сдаётся слишком рано и~теряет подписчиков.

\begin{figure}[tbp]
  \centering
  \includefiguresvg[width=.8\linewidth]{assets/figures/ch22-subscriptions/decline-decision-tree}
  \caption{Дерево решений при отказе MIT.}\label{fig:decline-decision-tree}
\end{figure}

\section{Обновление реквизитов}\label{sec:subscriptions-account-updater}

Сервисы обновления реквизитов закрывают сценарий, когда клиент не
менял подписку, но эмитент перевыпустил карту. Старый PAN больше
не подходит, а~продавец продолжает списывать по устаревшим данным.

Visa Account Updater (VAU) и Mastercard Automatic Billing Updater
(ABU) передают продавцу обновлённые реквизиты от участвующих
эмитентов, когда карта истекла или была
перевыпущена~\cite{mc-abu-privacy}. В~правилах Mastercard для
регулярных платежей участие в ABU формализовано, с~исключениями
по отдельным типам карт~\cite{mc-tpr}. Visa называет аналогичный
механизм credential lifecycle management. Обновление PAN и~сетевых
токенов возможно в~реальном времени прямо во время
авторизации~\cite{visa-lifecycle-management}.

Охват зависит от сети, рынка,
типа карты и~модели участия. Не каждый карточный продукт попадает в
один и~тот же механизм обновления, а~сами механизмы могут быть как
пакетными, так и встроенными прямо в~авторизацию~\cite{mc-tpr,visa-lifecycle-management}.

Разница между пакетным и авторизационным режимом существенна
для архитектуры. Пакетный режим, batch, работает по расписанию --
продавец или PSP отправляет файл с~реквизитами, схема возвращает
обновлённые данные через несколько часов. Для подписок с~месячным
циклом этого достаточно, если файл отправляется заблаговременно.
Авторизационный режим, real-time, встроен в~сам запрос. Если
реквизиты устарели, ответ содержит обновлённые данные, и~продавец
сохраняет их сразу. Visa описывает это как Real-Time Account
Updater~\cite{visa-lifecycle-management}; Mastercard --
как Real-Time ABU API; точное наименование сверяется
по действующей редакции TPR~\cite{mc-tpr}.

Покрытие зависит от участия эмитента. Не каждый банк подключён
к ABU или VAU, не каждый тип карты попадает в~процесс обновления.
Visa рекомендует подключать все BIN, включая reloadable prepaid,
но нерезагружаемые предоплаченные карты и~продукты малых эмитентов часто
остаются за пределами сервиса. Обновление реквизитов снижает
технический отток, но не устраняет его. Для карт вне покрытия
единственный путь -- попросить клиента ввести новые данные.

Сетевые токены меняют эту механику. Если продавец хранит не PAN,
а~токен схемы, при перевыпуске карты привязка обновляется
автоматически -- токен остаётся прежним, underlying PAN обновляется
на стороне схемы~\cite{visa-lifecycle-management}. Задача обновления
переносится от продавца к~схеме, а пакетные файлы для токенизированных
реквизитов становятся не нужны.

\begin{figure}[tbp]
  \centering
  \includefiguresvg[width=.8\linewidth]{assets/figures/ch22-subscriptions/credential-update-flow}
  \caption{Обновление реквизитов: VAU/ABU и~сетевые токены.}\label{fig:credential-update-flow}
\end{figure}

MIT уходит по карте с~истёкшим сроком, а~сервис
обновления реквизитов (ABU/VAU) недоступен -- пакетный файл
не обработался или реальновременный запрос вернул ошибку.
Эмитент отклоняет по коду~\texttt{54} (Expired Card).
Единственный путь -- уведомить клиента и попросить обновить
реквизиты вручную. Если продавец вместо этого повторяет тот же
запрос, он тратит попытки впустую и~рискует попасть
в~мониторинг схемы по избыточным ретраям.

Другой сценарий выглядит так. Сетевой токен обновился при перевыпуске
карты, underlying PAN изменился на стороне TSP, но первая же MIT
по обновлённому FPAN получает отказ с~кодом~\texttt{62}
(Restricted Card), потому что новая карта выпущена с~ограничениями
по MCC или по каналу CNP.
Продавец видит, что токен жив, но авторизация не проходит;
retry бесполезен, потому что ограничение на стороне эмитента.
Дальше нужно уведомить клиента, запросить новый платёжный
инструмент или, если доступно, обратиться к~эмитенту через
PSP для снятия ограничения.

Involuntary churn (отток из-за платёжных отказов, не по решению
клиента) по индустриальным оценкам Recurly, Spreedly и~других
вендоров биллинга обычно укладывается в 3--5\,\% подписочной
базы в~месяц; успешность ABU/VAU при обновлении реквизитов
по тем же оценкам -- порядка 85--90\,\% для карт участвующих
эмитентов. Это не публичные данные схем, а агрегированные
показатели по портфелям подписочных продуктов. Сетевые токены
повышают этот показатель, потому что обновление происходит
автоматически на стороне TSP без зависимости от пакетного цикла.

\section{Как выглядит устойчивый подписочный процесс}\label{sec:subscriptions-anti-churn}

Устойчивый процесс начинается с~первой CIT, и~она должна пройти,
корректно зафиксировав согласие. Система сохраняет идентификатор
исходной авторизации и токен реквизита, а~клиент получает ясное
описание условий подписки -- периодичность и сумму.

Каждый биллинговый цикл строится как MIT с~корректными флагами
и своевременной проверкой актуальности реквизитов. При перевыпуске
карты помогает ABU/VAU. Если эмитент просит вернуть клиента в~цикл
подтверждения, повторная попытка идёт по правилам схемы, а~не по
внутренней логике биллинга.

Корректная NTID-цепочка сигнализирует эмитенту, что запрос --
продолжение известного мандата. Эмитент применяет менее строгую
антифрод-политику, и~доля одобрений на MIT с~корректным NTID
выше, чем на \enquote{сиротских} запросах без истории.

Клиент должен заранее знать, когда будет следующая попытка списания
и~что ему делать, если платёж не прошёл. Понятное ожидание снижает
и~отток, и количество бессмысленных повторов.

\practicenote{%
  Устойчивость подписки держится на непрерывной цепочке согласия (CIT
  + идентификатор исходной авторизации) и актуальных реквизитах. Если
  одно из двух выпадает, эмитент теряет контекст операции.

}
